\begin{salyga}
Trikampio $ABC$ kraštinių $AB$ ir $AC$ vidurio taškai yra atitinkamai
$D$ ir $E$. Žinoma, kad $BC=10\,\text{cm}$.
Įrodykite, kad $DE\parallel BC$, ir apskaičiuokite atkarpos $DE$ ilgį.
\end{salyga}

\begin{atsakymas}
$DE\parallel BC$, o $DE=5\,\text{cm}$.
\end{atsakymas}

\begin{sprendimas}
Kadangi $D$ ir $E$ yra kraštinių vidurio taškai, tai
\[
  \frac{AD}{AB}=\frac{AE}{AC}=\frac12.
\]
Kampas $\angle DAE$ sutampa su kampu $\angle BAC$.
Todėl trikampiai $ADE$ ir $ABC$ yra panašūs pagal dvi proporcingas
kraštines ir kampą tarp jų. Viršūnes atitinka poros
$A\leftrightarrow A$, $D\leftrightarrow B$, $E\leftrightarrow C$.

Iš panašumo gauname $\angle ADE=\angle ABC$. Tiesė $AB$ kerta
tieses $DE$ ir $BC$, o šie kampai yra lygūs atitinkamieji kampai.
Vadinasi, $DE\parallel BC$.

Atitinkamų panašiųjų trikampių kraštinių santykiai lygūs, taigi
\[
  \frac{DE}{BC}=\frac{AD}{AB}=\frac12,
  \qquad DE=\frac{10}{2}=5\,\text{cm}.
\]
Taigi $DE\parallel BC$ ir $DE=5\,\text{cm}$.\hfill$\square$
\end{sprendimas}
